\documentclass[11pt]{article}
\usepackage[T1]{fontenc}
\usepackage{lmodern,amsmath,amssymb,geometry,microtype,xcolor,fancyhdr,hyperref}
\geometry{letterpaper,left=0.96in,right=0.96in,top=0.92in,bottom=0.86in}
\definecolor{ink}{HTML}{183953}\definecolor{tone}{HTML}{386475}
\hypersetup{colorlinks=true,urlcolor=tone,pdftitle={A matching formula for Schur products and a weighted Chollet inequality}}
\setlength{\parindent}{0pt}\setlength{\parskip}{5pt}
\pagestyle{fancy}\fancyhf{}\setlength{\headheight}{15pt}
\lhead{\small\sffamily\color{tone} MATRIX PERMANENTS}\rhead{\small\sffamily\color{tone} RESEARCH NOTE}
\rfoot{\small\sffamily\color{tone}\thepage\,/\,3}
\newcommand{\shead}[1]{\par\vspace{4pt}{\large\bfseries\sffamily\color{ink}#1}\par}
\newcommand{\per}{\operatorname{per}}
\newcommand{\dd}{\mathcal D_n}
\begin{document}
{\small\sffamily\color{tone}18 NOVEMBER 2026 \quad / \quad RESEARCH NOTE}\par
{\LARGE\bfseries\sffamily\color{ink} A matching formula for Schur products\\[2pt] and a weighted Chollet inequality}\par
\vspace{5pt}
For a fixed correlation matrix, what is the largest permanent obtainable by taking its entrywise product with a scaled diagonally dominant correlation matrix? It is the exponential of the maximum matching weight, with edge weights $\log(1+|a_{ij}|^2)$. In particular, every weighted graph Laplacian satisfies the strong self-Chollet inequality.
\shead{1. Exact extremal statement}
A correlation matrix is a complex Hermitian positive semidefinite matrix with diagonal entries one. Write $A\circ B=(a_{ij}b_{ij})$, $|A|=(|a_{ij}|)$, and $\per A=\sum_\pi\prod_i a_{i,\pi(i)}$. Define
\[
  \dd=\{B\text{ a correlation matrix}:\ \rho(C)\le 1,
  \quad c_{ii}=0,\quad c_{ij}=|b_{ij}|\ (i\ne j)\}.
\]
Here $\rho$ is the spectral radius. For symmetric, nonnegative $C$, this is exactly the condition $I-C\succeq0$ on the comparison matrix. A matching $M$ is any set of mutually vertex-disjoint unordered pairs in $\{1,\ldots,n\}$; the empty product is one.
  
\textbf{Theorem.} For every correlation matrix $A$,
\begin{equation}
\max_{B\in\dd}\per(A\circ B)
=\max_{B\in\dd}\per\bigl(|A\circ B|\bigr)
=\max_M\prod_{\{i,j\}\in M}(1+|a_{ij}|^2)
\ \le\ \per A.
\end{equation}
A maximizing $B$ consists of $2\times2$ all-ones blocks on a maximizing matching and singleton identity blocks elsewhere.

\textbf{Graph consequence.} Give a finite undirected loopless graph arbitrary real edge weights $w_{ij}=w_{ji}\ge0$. Set $d_i=\sum_j w_{ij}$, $L_{ii}=d_i$, and $L_{ij}=-w_{ij}$ off the diagonal. For $H=L+\operatorname{diag}(\alpha_i)$, $\alpha_i\ge0$, and every vertex set $S$,
\begin{equation}
 \per(H[S]\circ H[S])\le \per(H[S])\prod_{i\in S}(d_i+\alpha_i)
 \le \per(H[S])^2.
\end{equation}
The diagonal entries retain contributions from edges leaving $S$. This includes arbitrary weighted nonbipartite blocks and singular Laplacians. In fact (2), with the actual diagonal entries on the right, holds for every real symmetric positive semidefinite matrix with nonpositive off-diagonal entries.

\textbf{Scope.} Pant and Singh ask the unweighted graph question in [1, \S6]. The repository [2] treats that case and separately proves weighted cycle/bipartite block families; it explicitly leaves general weighted graphs and positive semidefinite $Z$-matrices unresolved in that work. Equation (2) covers those classes. It is not the unrestricted Chollet conjecture for two arbitrary positive semidefinite matrices.
\newpage
\shead{2. From random arrows to a matching}
Fix $A$ and $B\in\dd$. On each connected component of the support of $C$ containing an edge, the Perron--Frobenius theorem gives a strictly positive eigenvector of that component with eigenvalue its spectral radius, which is at most one. Take value one on each isolated vertex. Combining the vectors gives $v>0$ with $Cv\le v$, including when the global radius equals one or component radii differ.

Independently for each vertex $i$, draw one outgoing arrow $i\to j$ with probability
\[
 q_{ij}=c_{ij}\frac{v_j}{v_i}\quad(i\ne j).
\]
With the remaining probability the vertex has no outgoing arrow. These are valid probabilities by $Cv\le v$. Call the random directed graph $F$. Its directed cycles are disjoint, have length at least two, and are counted up to rotation; a two-cycle is counted once. For a cycle $D$, set $u(D)=\prod_{i\to j\in D}|a_{ij}|$.

\textbf{Cycle identity.} The following equality is exact, including zeros:
\begin{equation}
\per(|A\circ B|)=
 \mathbb E\!\left[\prod_{D\ \mathrm{cycle\ of}\ F}(1+u(D))\right].
\end{equation}
Indeed, each permutation is a set of vertex-disjoint directed cycles of length at least two, with every other vertex fixed. Expand the product on the right into its cycle subsets. For any fixed such family, the probability that all its required arrows are present is the product of their $q_{ij}$, by row independence. The factors $v_j/v_i$ telescope on every cycle, so this probability equals the product of the $|b_{ij}|$ on the family. Multiplication by the $u(D)$ yields exactly its permanent monomial; the fixed vertices contribute one. Extra arrows, and extra cycles not chosen in the expansion, impose no restrictions. Summing proves (3).

\textbf{Pointwise bound.} On each cycle $D$ of a realized $F$, select an unordered edge $e(D)=\{i,j\}$ maximizing $|a_{ij}|$ there. A correlation matrix has $|a_{ij}|\le1$ by its $2\times2$ principal minors. Thus for $k=|D|\ge2$ and $t=|a_{ij}|$ on the chosen edge, $u(D)\le t^k\le t^2$. Since the cycles are disjoint, their selected edges form a matching. Therefore, for every realization,
\[
 \prod_D(1+u(D))\le \prod_D(1+|a_{e(D)}|^2)
 \le \max_M\prod_{\{i,j\}\in M}(1+|a_{ij}|^2).
\]
Together with (3) and the triangle inequality, this bounds both permanents in (1) from above. Conversely, choose $B$ with the blocks specified after (1). The comparison graph is a matching, so $\rho(C)\le1$, and these blocks are positive semidefinite. Both permanents factor over the same blocks, attaining exactly the claimed value.

Finally, Lieb's proved block permanent inequality [3] gives $\per A\ge\prod_j\per A[S_j]$ for a partition into principal blocks: iterate its two-block form. Use the matched pairs as blocks and the unmatched vertices as singletons. The pair permanent is $1+|a_{ij}|^2$, and each singleton contributes one. This proves the last inequality in (1).
\newpage
\shead{3. Undoing the normalization}
Let $X,Y$ be Hermitian positive semidefinite matrices of the same size. If either has a zero diagonal entry then its entire corresponding row is zero by its $2\times2$ minors. Both sides of $\per(X\circ Y)\le\per(X)\prod_i y_{ii}$ then vanish. Otherwise put
\[
 a_{ij}=\frac{x_{ij}}{\sqrt{x_{ii}x_{jj}}},\qquad
 b_{ij}=\frac{y_{ij}}{\sqrt{y_{ii}y_{jj}}}.
\]
If the resulting $B$ satisfies the comparison condition in (1), then positive diagonal scaling in the permanent gives
\begin{equation}
 \per(X\circ Y)\le\per(X)\prod_i y_{ii}\le\per(X)\per(Y).
\end{equation}
The second inequality follows from the singleton case of Lieb's block inequality. No comparison condition is imposed on $X$. For a principal submatrix, the normalized $C$ is restricted to that principal set; $I-C\succeq0$ is preserved.

If $Y$ is a real symmetric positive semidefinite $Z$-matrix (off-diagonal entries nonpositive), its normalization equals $I-C$. Consequently $\lambda_{\max}(C)\le1$, and for nonnegative symmetric $C$ this equals $\rho(C)$. Equation (4) applies. Weighted graph Laplacians and their nonnegative diagonal additions are of this type, since $z^*Lz=\sum_{i<j} w_{ij}|z_i-z_j|^2$. Taking $X=Y=H[S]$ proves (2). The empty matrix convention gives one on all sides.

\textbf{A weighted nonbipartite example.} On $K_4$, give edges $12,13,14,23,24,34$ weights $1,2,3,4,5,6$ respectively. The degree vector is $(6,10,12,14)$. Direct expansion gives $\per L=15{,}232$ and
\[
 \per(L\circ L)=114{,}290{,}848
 \ <\ (6\cdot10\cdot12\cdot14)\per L=153{,}538{,}560 .
\]
This is an illustration, not evidence needed for the universal proof.
\shead{Sources and limits}
{\small
[1] P. Pant and R. Singh, \href{https://arxiv.org/html/2604.24192v2}{\emph{Structural Classes for Chollet's Permanent Conjecture}}, arXiv:2604.24192v2 (2026), especially Lemma 1 and \S6.

[2] Y. Zhang, \href{https://github.com/mxym/math/tree/main/preprints/all-graph-chollet-2026-10}{\emph{Strong Chollet inequalities for all simple graph Laplacians}} and \href{https://github.com/mxym/math/blob/main/research/weighted-cycle-chollet/PAPER.md}{\emph{Strong Chollet inequalities for arbitrarily weighted cycles and block graphs}} (2026), especially \S7 of the latter. The latter uses independent neighbor choices for mutual-edge matchings; it records a numerical matching-bound candidate, not an all-weight proof.

[3] E. H. Lieb, \href{https://doi.org/10.1512/iumj.1967.16.16008}{\emph{Proofs of some conjectures on permanents}}, J. Math. Mech. 16 (1966), 127--134. The needed two-block statement is also explicit in [1, Lemma 1].

[4] S. W. Drury, \href{https://doi.org/10.13001/1081-3810.3178}{\emph{A counterexample to a question of Bapat and Sunder}}, Electron. J. Linear Algebra 31 (2016), 69--70. Without the comparison restriction the first inequality of (4) is false in general. The weaker unrestricted Chollet inequality is not settled by this argument.

[5] K. H. Kim and F. W. Roush, \href{https://doi.org/10.1016/0024-3795(81)90090-2}{\emph{Expressions for certain minors and permanents}} (1981); Z. Chen and L. Cao, \href{https://arxiv.org/pdf/1801.00225}{\emph{On the maximum of the permanent of $(I-A)$}} (2018), Lemma 3.5. Random-map expansions and cycle factorization on digraphs with at most one outgoing edge per vertex are prior techniques. Here the second matrix supplies the arrow probabilities; the first supplies the cycle weights and the matching objective.

This AI-assisted research note supplies a written proof, not external peer review, formal verification, or a historical priority claim.}
\end{document}
